\section{Conclusión}

El diseño de las losas del piso tipo y de la viga del subterráneo cumple con los requisitos de resistencia y deformación establecidos.

En cuanto a deformaciones, la losa 302 en niveles con hormigón G25 resulta ser la más desfavorable, superando levemente el límite admisible de deformación, por lo que se requiere incorporar una contraflecha reducida. El resto de las losas cumple sin necesidad de medidas adicionales. Para la viga continua, las deformaciones calculadas se mantienen dentro de los límites admisibles en ambas condiciones de apoyo evaluadas.

La armadura mínima de $A_{s,\text{mín}} = 2{,}70$~cm²/m gobierna en la totalidad de las losas 305i y 305ii, y en la dirección $x$ positiva de varios paños intermedios. Para la viga continua, la armadura mínima queda ampliamente superada por la armadura adoptada de $4\phi25$.

Las armaduras se especifican utilizando un diámetro único $\phi$8, con separaciones entre 10 y 19~cm. La losa más exigente corresponde a la 302, con $\phi$8~@~10~cm en momentos negativos para los niveles con hormigón G25. La viga se diseña con $4\phi25$ a flexión y $\phi10$~@~225~mm a corte.

El espesor de 15~cm adoptado para el piso tipo resulta adecuado. La magnitud de la contraflecha requerida en la losa 302 es baja, por lo que no resulta necesario modificar el espesor adoptado, verificándose además un comportamiento adecuado en el resto de los paños.